\section{Materials and Methods}

\subsection{Database and record-level partition}
Experiments used the MIT-BIH Arrhythmia Database distributed through PhysioNet, which contains 48 half-hour, two-lead ambulatory ECG recordings sampled at 360~Hz\cite{goldberger2000physiobank,moody2001impact}. Following the AAMI-oriented preprocessing implemented in the repository, paced-beat recordings were excluded and 45 records were retained for model development and evaluation. A fixed record-level split (\texttt{test\_size}=0.15; random seed 42) assigned seven records (106, 207, 208, 220, 228, 231, and 233) to the held-out test cohort and the remaining 38 records to the development pool. The held-out records contained 15{,}573 beats and were excluded from hyperparameter selection and model fitting. Five-fold cross-validation was performed within the 38-record development pool, after which the selected configuration was fitted on the development records and evaluated on the held-out cohort. Table~\ref{tab:cohort-distribution} summarises the evaluation cohorts and the final test-set class support.

The record list was fixed before tuning and is read by the hyperparameter-tuning, cross-validation, and final-evaluation code. Data acquisition is handled by \texttt{download\_data.py}; the split identifiers and run artifacts are retained with the repository.

\subsection{Beat representation and sequence construction}
The preprocessing pipeline (\texttt{preprocess\_data.py}) detects R peaks, extracts a 187-sample window centred on each annotated beat from the first ECG channel (approximately 520~ms at 360~Hz), and computes a continuous wavelet transform using the PyWavelets Morlet kernel (\texttt{morl}) at 32 integer scales. Beat localisation follows established QRS-detection principles, and the wavelet transform describes transient morphology across scales\cite{pan1985real,laguna1994automatic,addison2005wavelet,torrence1998practical,mallat1989wavelet}. The resulting time--scale map has dimensions $32\times187$.

Cached scalograms remain unnormalised. During dataset construction, amplitude normalisation uses statistics from the relevant development data. The TFRecord writer (\texttt{create\_batched\_tfrecords.py}) groups three temporally adjacent beats and stores the tensor shape and sequence metadata needed for streaming. This short sequence gives the classifier immediate rhythm context without greatly extending the input.

\subsection{CNN--Conformer and comparator models}
The primary classifier is implemented in \texttt{ModelBuilder.py}. Two-dimensional convolutional blocks encode local patterns in each Morlet map. Conformer blocks then process the three-beat representation using macaron-style feed-forward modules, multi-head self-attention, depthwise convolution, residual connections, layer normalisation, and dropout\cite{gulati2020conformer,vaswani2017attention,he2016deep,srivastava2014dropout,ba2016layer}. A dense output layer gives probabilities for the five AAMI categories N, S, V, F, and Q.

Three comparators were evaluated on the same held-out records. They comprised an attention-only neural network, a CNN--LSTM in which recurrent modelling replaces the Conformer sequence encoder, and a feature-engineered classifier. The comparison is confined to these four implementations and was designed to assess the hybrid model within a common pipeline, not to reproduce the wider published literature.

\subsection{Optimisation and model selection}
Hyperparameter search used Hyperband\cite{li2018hyperband}. The search covered embedding dimension, number of attention heads, convolution kernel width, dropout, and the number of Conformer blocks. \texttt{DataLoader.py} streamed TFRecord batches and applied fold-specific normalisation. Neural models used AdamW at a learning rate of $1\times10^{-5}$, sparse categorical cross-entropy, and class weights calculated from the development labels. Early stopping and learning-rate scheduling were used during the Hyperband trials. For the definitive evaluation, the selected configuration was trained for 20 epochs on the development data without early stopping or gradient clipping. Mixed-float16 execution was enabled on compatible GPUs. The random seed, training history, fitted model, classification report, and raw test probabilities were archived under \path{Research_Runs}.

\subsection{Evaluation and uncertainty}
Held-out performance was described by accuracy, macro-F1, weighted-F1, class-wise precision, recall and F1, and one-vs-rest area under the receiver operating characteristic curve (AUC). Macro-F1 assigns equal weight to each rhythm class and is less affected by the abundance of normal beats than accuracy or weighted-F1. Accuracy across the five development folds is reported as mean and standard deviation. Where archived with the study tables, held-out point estimates are accompanied by 95\% bootstrap confidence intervals. Receiver operating characteristic and precision--recall curves were generated from the stored probabilities; the latter are especially informative for the sparse classes\cite{fawcett2006introduction,saito2015precision}.

Two exploratory analyses were performed after the definitive predictions had been generated. One-vs-rest thresholds for SVEB and fusion beats were swept on the stored test probabilities to examine whether the default operating point accounted for their low F1-scores. Per-class expected calibration error (ECE; 15 equally spaced confidence bins) and Brier scores were also calculated, followed by single-parameter temperature scaling. Prediction certainty and calibrated uncertainty have been treated as explicit outputs in recent AI-enabled ECG studies\cite{demolder2024certainty,doggart2025singlelead}. Here, both analyses are post-hoc: the threshold sweep used test predictions, and the temperature was fitted on the same test cohort. The results describe this dataset and do not constitute independent threshold selection or calibration validation.

\subsection{Figure preparation}
OpenAI ChatGPT (GPT-5.6 Sol) assisted with the LaTeX/TikZ layout of the workflow and summary figures; it did not generate or modify data, predictions, or reported values. Values in those figures were transcribed from the archived outputs listed in \texttt{qa/claim-ledger.md}. Supplementary confusion-matrix, receiver operating characteristic, and precision--recall plots were regenerated deterministically from the archived held-out probabilities. No graphical abstract generated by a general-purpose AI tool is included.

\begin{figure*}[t]
  \centering
  \resizebox{0.94\textwidth}{!}{% AI-assisted explanatory figure. Counts are direct transcriptions from the study protocol.
\begin{tikzpicture}[
  node distance=0.9cm and 1.15cm,
  base/.style={rectangle,rounded corners=2.5pt,minimum height=1.25cm,minimum width=2.75cm,text width=2.5cm,align=center,font=\footnotesize,thick},
  source/.style={base,draw=black!65,fill=black!3},
  dev/.style={base,draw=blue!70!black,fill=blue!6},
  test/.style={base,draw=green!55!black,fill=green!7},
  diag/.style={base,draw=black!45,fill=black!2,dashed},
  arrow/.style={-{Stealth[length=4pt]},thick,draw=black!72}
]

\node[source] (source) {\textbf{MIT-BIH}\\[2pt]{\Large 45} retained\\non-paced records};
\node[source,right=of source] (split) {\textbf{Fixed record split}\\seed 42};

\node[dev,right=1.4cm of split,yshift=1.35cm] (devpool) {\textbf{Development}\\[2pt]{\Large 38} records};
\node[dev,right=of devpool] (repr) {\textbf{Representation}\\187-sample beats\\32-scale Morlet\\3-beat context};
\node[dev,right=of repr] (fit) {\textbf{Model development}\\Hyperband\\5-fold CV\\final fitting};

\node[test,right=1.4cm of split,yshift=-1.35cm] (holdout) {\textbf{Held-out test}\\[2pt]{\Large 7} records\\{\large 15{,}573} beats};
\node[test,right=of holdout] (eval) {\textbf{One final evaluation}\\accuracy\\macro/weighted F1\\class AUC};
\node[test,right=of eval] (finding) {\textbf{Primary result}\\baseline comparison\\+\\rare-rhythm failures};
\node[diag,below=0.85cm of finding] (posthoc) {\textbf{Post-hoc diagnostics}\\threshold sweeps\\ECE / Brier};

\draw[arrow] (source) -- (split);
\draw[arrow] (split.east) -- ++(0.55,0) |- (devpool.west);
\draw[arrow] (split.east) -- ++(0.55,0) |- (holdout.west);
\draw[arrow] (devpool) -- (repr);
\draw[arrow] (repr) -- (fit);
\draw[arrow] (holdout) -- (eval);
\draw[arrow] (fit.south) |- (eval.north);
\draw[arrow] (eval) -- (finding);
\draw[arrow,dashed] (finding) -- (posthoc);

\node[anchor=east,font=\small\bfseries,text=blue!70!black] at ($(devpool.west)+(-0.22,0)$) {DEVELOPMENT};
\node[anchor=east,font=\small\bfseries,text=green!45!black] at ($(holdout.west)+(-0.22,0)$) {HELD-OUT};
\end{tikzpicture}
}
  \caption{Study workflow. Thirty-eight records form the development pool, and seven records remain separate until final evaluation. Threshold and calibration analyses use the stored test predictions and are post-hoc. The TikZ layout was AI-assisted; counts come from the archived study protocol.}
  \label{fig:pipeline}
\end{figure*}
